\section{Problem Evolution}
\label{sec:problem-evolution}

The intent of this section is to follow the path of dynamics evolution.
Understanding that path makes the design issues and nomenclature visible.  It
also helps future development because the current model contains assumptions
about rules such as equal-radius particles, wall particles, and source-owned
contact state.  When those rules change, it should be clear where they came
from and why they existed.

\subsection{Observed Failure - Velocity Asymmetry After Collision in TwoParticleHorizontal}

\noindent\textbf{Commit hash:} \texttt{62ff53dbe45b11fb5aad8f5695d5f9513d2baf2a}

The first Geo collision response used a purely frame-local overlap impulse:

\begin{equation}
J = q A \Delta t.
\end{equation}

Here $q$ is the particle-owned pair stiffness, $A$ is the current overlap area,
and $\Delta t$ is the frame time step.  This rule was intentionally simple and
source-owned: each shader-shaped invocation updated only the current source
particle.

In \texttt{TwoParticleHorizontal}, total momentum remained conserved because
the two particles received equal and opposite velocity changes.  However, the
final speeds drifted slightly high:

\begin{equation}
v = 0.05 \quad \rightarrow \quad v \approx 0.050006.
\end{equation}

This is symmetric velocity drift.  It is invisible to total momentum:

\begin{equation}
P = \sum_i m_i v_i \approx 0,
\end{equation}

but visible in kinetic energy:

\begin{equation}
K = \sum_i \frac{1}{2}m_i \lVert v_i \rVert^2.
\end{equation}

The drift appeared at the end of rebound.  The particles were already
separating, but they were still slightly overlapped.  Because the old rule
continued applying $qA\Delta t$ until the overlap disappeared, the last small
overlap frames released a little extra velocity.

\subsection{Rejected Pairwise Fix}

A conventional pairwise elastic collision impulse can conserve both momentum
and kinetic energy:

\begin{equation}
J =
-\frac{(1+e)v_n}{\frac{1}{m_s}+\frac{1}{m_t}}.
\end{equation}

That approach updates both particles in one pair solve.  It is not allowed in
the current Geo path because the GLSL execution model is source-owned: one
invocation owns one source particle and must not write the target particle.

Therefore Python may not use a pairwise write-back shortcut.  The Python Geo
path must not become a separate reference model that fixes dynamics in a way
the shader cannot reproduce safely.

\subsection{Allowed Exception}

The accepted fix is source-owned internal normal momentum.  Each source-particle
contact slot may carry one persistent scalar:

\begin{equation}
I_{s,t} \ge 0.
\end{equation}

This value records the normal momentum removed from the source during
compression.  It is not a shared pair reservoir and it does not write the
target particle.

In the GLSL-shaped contact state, the value is stored in:

\begin{verbatim}
P[SourceID].contacts[slot].aux.z
\end{verbatim}

\subsection{Compression Rule}

The raw geometric impulse remains:

\begin{equation}
J_{raw} = q_{s,t} A_{s,t} \Delta t.
\end{equation}

During compression, the source can store only momentum that actually exists in
the source normal motion:

\begin{equation}
J_c =
\min\left(
J_{raw},
\max(0, m_s(v_s \cdot \hat{n}_{s,t}))
\right).
\end{equation}

The source velocity is updated by:

\begin{equation}
v_s' = v_s - \frac{J_c}{m_s}\hat{n}_{s,t}.
\end{equation}

The source-owned internal momentum ledger increases by the same amount:

\begin{equation}
I_{s,t}' = I_{s,t} + J_c.
\end{equation}

\subsection{Returning Rule}

When the source can no longer provide the full compression impulse, the contact
enters the returning phase.  During returning, the source may release only
stored internal momentum:

\begin{equation}
J_r = \min(J_{raw}, I_{s,t}).
\end{equation}

The velocity update uses the same source-owned write shape:

\begin{equation}
v_s' = v_s - \frac{J_r}{m_s}\hat{n}_{s,t}.
\end{equation}

The ledger is reduced:

\begin{equation}
I_{s,t}' = \max(0, I_{s,t} - J_r).
\end{equation}

This cap prevents the last tiny-overlap frame from releasing more momentum than
it stored during compression.

\subsection{Result}

For \texttt{TwoParticleHorizontal}, the internal momentum ledger drains to zero
at the end of contact and the kinetic energy returns to the starting value:

\begin{verbatim}
start_ke=0.002500000000000000
final_ke=0.002500000000000000
ke_drift=-8.67e-19
p1 vx=-0.050000000000000
p2 vx=+0.050000000000000
\end{verbatim}

The fix preserves the current GLSL architecture:

\begin{itemize}[leftmargin=*]
    \item each invocation still owns only one source particle;
    \item the target particle is read-only;
    \item the carried value is source-owned contact state;
    \item rebound cannot release more momentum than compression stored.
\end{itemize}
